\section{Закон инерции Сильвестра}

\begin{remark}
    Напомним: квадратичная форма $Q$ над $\mathbb{R}$ в некотором базисе приводится к виду
    \[
        Q = x_1^2 + \dots + x_p^2 - x_{p+1}^2 - \dots - x_{p+m}^2.
    \]
    Пара $(p, m)$ называется \emph{сигнатурой} $Q$.
\end{remark}


\begin{theorem}[Закон инерции Сильвестра]
    Сигнатура $(p, m)$ не зависит от выбора базиса.

    \begin{proof}
        \textbf{Шаг 1: сведение к невырожденному случаю.}
        Рассмотрим $\ker Q$ (ядро ассоциированной билинейной формы) и форму $\bar{Q}$ на $V / \ker Q$.
        Она невырожденна. Сигнатура $\bar{Q}$ определяет сигнатуру $Q$ (добавляются нули).
        НУО $Q$ невырожденна, и $p + m = \dim V$.

        \textbf{Шаг 2: разложение в гиперболическую и анизотропную части.}
        НУО $p \geq m$ (иначе рассмотрим $-Q$). Сгруппируем слагаемые попарно:
        \[
            Q = (x_1^2 - x_{p+1}^2) + (x_2^2 - x_{p+2}^2) + \dots + (x_m^2 - x_{p+m}^2) + x_{m+1}^2 + \dots + x_p^2.
        \]

        Для каждой пары $(e_i, e_{p+i})$ рассмотрим плоскость $H_i = \langle e_i, e_{p+i} \rangle$.
        Матрица Грама $Q\big|_{H_i}$ принимает вид $\bigl(\begin{smallmatrix} 0 & 1 \\ 1 & 0 \end{smallmatrix}\bigr)$ --- гиперболическая форма.
        В базисе $f_i = \frac{e_i + e_{p+i}}{\sqrt{2}}$: \\
        \[
            \left(
            \begin{array}{ccc|c}
            \fbox{$\begin{smallmatrix} 0 & 1 \\ 1 & 0 \end{smallmatrix}$} & & & \\
            & \fbox{$\begin{smallmatrix} 0 & 1 \\ 1 & 0 \end{smallmatrix}$} & & \text{\huge 0} \\
            & & \ddots & \\ \hline
            & 0 & & E
            \end{array}
            \right)
        \]
        , в $\tilde{f}_i = \frac{e_i - e_{p+i}}{\sqrt{2}}$:

        \[
            \left(
            \begin{array}{c|c}
            \begin{array}{c|c}
            0 & E \\ \hline
            E & 0
            \end{array} & 0 \\ \hline
            \text{0} & E
            \end{array}
            \right)
        \]


        Подпространство $A = \langle e_{m+1}, \dots, e_p \rangle$ анизотропно:
        для $v = \sum \alpha_i e_i$ имеем $Q(v) = \sum \alpha_i^2 > 0$ при $v \neq 0$.

        Итого $V = H \oplus A$, где $H = H_1 \oplus \dots \oplus H_m$ гиперболическое, $A$ анизотропное.

        \textbf{Шаг 3: единственность разложения.} Из леммы Витта следует,
        что разложение $V = H \oplus A$ единственно с точностью до изометрии.
        Действительно, пусть $V = \tilde{H} \oplus \tilde{A}$ --- другое такое разложение.
        Гиперболические подпространства изометричны тогда и только тогда, когда у них совпадают размерности.
        Если $\dim H > \dim \tilde{H}$, выделим в $H$ подпространство $\tilde{H} \oplus H'$ и по лемме Витта получим
        $H' \oplus A \cong \tilde{A}$. Но в $H'$ есть изотропный вектор, а в $\tilde{A}$ нет --- противоречие.
        Значит $\dim H = \dim \tilde{H}$, и по лемме Витта $A \cong \tilde{A}$.

        Поэтому $m = \dim H / 2$ и $p = \dim A + m$ определяются однозначно.
    \end{proof}
\end{theorem}
